\documentclass[10pt,a4paper]{article}

% ==============================================================================
% PAQUETES DE CODIFICACIÓN Y TIPOGRAFÍA
% ==============================================================================
\usepackage{fontspec}
\setmainfont{HelveticaNeue}[
  Path = ./fonts/,
  Extension = .otf,
  UprightFont = *-Regular,
  BoldFont = *-Bold,
  ItalicFont = *-RegularItalic
]

\usepackage[spanish,es-noshorthands]{babel}
\usepackage[a4paper, left=1.5cm, right=1.5cm, top=1.5cm, bottom=1.5cm]{geometry}
\usepackage{xcolor}
\usepackage[hidelinks]{hyperref}

\pagestyle{empty}
\setlength{\parindent}{0pt}
\setlength{\parskip}{0pt}
\linespread{1.16}

% ==============================================================================
% 1. CONFIGURACIÓN DE COLORES
% (Puedes cambiar los códigos hexadecimales aquí si deseas ajustar la paleta)
% ==============================================================================
\definecolor{accent}{HTML}{0051A3}       % Azul corporativo (títulos de sección, nombres de empresas)
\definecolor{dotaccent}{HTML}{0071E3}    % Azul vibrante (punto del nombre, enlaces)
\definecolor{textdark}{HTML}{1D1D1F}     % Grafito oscuro principal (fidelidad Apple)
\definecolor{subtext}{HTML}{46464E}      % Gris intermedio (cargo, subtítulo)
\definecolor{subsecondary}{HTML}{6E6E73} % Gris secundario (fechas, modalidades, separadores)
\definecolor{linecolor}{HTML}{DCDFE4}    % Gris tenue (líneas divisorias de secciones)

% ==============================================================================
% 2. CONFIGURACIÓN DE ESPACIOS (Ajusta estos valores para cambiar las distancias)
% ==============================================================================
\newlength{\SpaceAboveSection}
\setlength{\SpaceAboveSection}{0.48cm}         % Espacio antes de cada sección

\newlength{\SpaceSectionTextToLine}
\setlength{\SpaceSectionTextToLine}{0.28cm}    % Separación equilibrada entre el título y la línea divisoria

\newlength{\SpaceSectionLineToContent}
\setlength{\SpaceSectionLineToContent}{0.28cm} % Separación entre la línea divisoria y el contenido

\newlength{\SpaceBetweenJobs}
\setlength{\SpaceBetweenJobs}{0.34cm}          % Espacio entre diferentes empleos/experiencias

\newlength{\SpaceJobTitleToSub}
\setlength{\SpaceJobTitleToSub}{0.08cm}        % Espacio entre cargo/empresa y ubicación/modalidad

\newlength{\SpaceJobSubToItems}
\setlength{\SpaceJobSubToItems}{0.20cm}        % Espacio entre ubicación/modalidad y la primera viñeta

\newlength{\SpaceBetweenItems}
\setlength{\SpaceBetweenItems}{0.20cm}         % Espacio vertical entre viñetas individuales

% ==============================================================================
% 3. DATOS PERSONALES (Modifica fácilmente tus datos aquí)
% ==============================================================================
\newcommand{\cvName}{Edwyn Guzmán}
\newcommand{\cvRole}{Desarrollador full-stack y DevOps \textbar\ React.js · Laravel · React Native · Docker · GitHub Actions}
\newcommand{\cvLocation}{Caracas, Venezuela (trabajo remoto)}
\newcommand{\cvPhone}{+58 424 154 3777}
\newcommand{\cvPhoneLink}{+584241543777}
\newcommand{\cvEmail}{edwynjesus30@gmail.com}
\newcommand{\cvLinkedInUser}{linkedin.com/in/edwyn-guzman-95324632b}
\newcommand{\cvLinkedInLink}{https://linkedin.com/in/edwyn-guzman-95324632b}
\newcommand{\cvGitHubUser}{github.com/edwynG}
\newcommand{\cvGitHubLink}{https://github.com/edwynG}

% ==============================================================================
% 4. MACROS REUTILIZABLES (Comandos para estructurar el contenido)
% ==============================================================================
\newcommand{\cvdot}{\textcolor{subsecondary}{~·~}}

% Encabezado del CV
\newcommand{\makecvheader}{%
  {\fontsize{24pt}{29pt}\bfseries \cvName{\color{dotaccent}.}}\par\vspace{0.26cm}%
  {\fontsize{10.2pt}{14.5pt}\color{subtext}\selectfont \cvRole}\par\vspace{0.22cm}%
  {\fontsize{9.4pt}{13.5pt}\color{subsecondary}\selectfont \cvLocation\cvdot\href{tel:\cvPhoneLink}{\color{subsecondary}\cvPhone}\cvdot\href{mailto:\cvEmail}{\color{dotaccent}\cvEmail}}\par\vspace{0.10cm}%
  {\fontsize{9.4pt}{13.5pt}\color{subsecondary}\selectfont LinkedIn: \href{\cvLinkedInLink}{\color{subsecondary}\cvLinkedInUser}\cvdot GitHub: \href{\cvGitHubLink}{\color{subsecondary}\cvGitHubUser}}\par%
}

% Encabezado de sección con tracking y línea divisoria con separación equilibrada
\newcommand{\cvsection}[1]{%
  \vspace{\SpaceAboveSection}%
  {\color{accent}\bfseries\fontsize{10.2pt}{12pt}\addfontfeatures{LetterSpace=9.0}\selectfont #1}\par\nopagebreak%
  \vspace{\SpaceSectionTextToLine}%
  {\color{linecolor}\hrule height 0.65pt width \linewidth}\par\nopagebreak%
  \vspace{\SpaceSectionLineToContent}%
}

% Entrada de trabajo o educación:
% Uso: \cventry{Cargo}{Empresa o Institución}{Ubicación · Modalidad}{Fechas}
\newcommand{\cventry}[4]{%
  {\fontsize{9.8pt}{13.5pt}\selectfont{\bfseries #1}\cvdot{\color{accent}\bfseries #2} \hfill {\color{subsecondary}\small #4}}\par%
  \vspace{\SpaceJobTitleToSub}%
  {\color{subsecondary}\small #3}\par%
  \vspace{\SpaceJobSubToItems}%
}

% Viñeta individual de logro o responsabilidad
% Uso: \cvitem{Texto del logro...}
\newcommand{\cvitem}[1]{%
  \par\hangindent=1.2em\hangafter=1\noindent
  \hbox to 1.2em{\raisebox{0.05ex}{\small\textbullet}\hss}%
  {\fontsize{9.6pt}{14.5pt}\selectfont #1}\par%
  \vspace{\SpaceBetweenItems}%
}

% Separación entre bloques de empleo
\newcommand{\jobsep}{\vspace{\SpaceBetweenJobs}}

% Fila de habilidades técnicas
% Uso: \cvskill{Categoría}{Herramientas...}
\newcommand{\cvskill}[2]{%
  {\fontsize{9.6pt}{15.2pt}\selectfont \textbf{#1:} #2}\par\vspace{0.14cm}%
}

% ==============================================================================
% CUERPO DEL DOCUMENTO
% ==============================================================================
\begin{document}
\color{textdark}

% ==============================================================================
% PÁGINA 1
% ==============================================================================

\makecvheader

% --- RESUMEN PROFESIONAL ---
\cvsection{RESUMEN PROFESIONAL}

{\fontsize{9.6pt}{14.8pt}\selectfont
  Desarrollador full-stack y DevOps, estudiante de licenciatura en computación en la Universidad Central de Venezuela y preparador universitario. Experiencia construyendo aplicaciones web y móviles en producción, diseño de APIs robustas y automatización de infraestructura. Especializado en React.js, Next.js, React Native, y Laravel, con sólida práctica en containerización con Docker, despliegue continuo mediante GitHub Actions y migraciones de sistemas hacia servidores locales y VPS. Base matemática y algorítmica aplicada a la optimización de código, gestión eficiente de bases de datos y fundamentos de aprendizaje automático.
}\par

% --- EXPERIENCIA PROFESIONAL ---
\cvsection{EXPERIENCIA PROFESIONAL}

\cventry{Desarrollador full-stack y DevOps}{GanaderaSoft}{Venezuela · Remoto · Proyecto de investigación UCV · FONACIT / MINCYT}{Noviembre 2025 – Septiembre 2026}
\cvitem{Reingeniería integral de todo el sistema, rediseñando la arquitectura para modernizar la plataforma y optimizar la automatización de datos productivos en ganadería vacuna.}
\cvitem{Reconstrucción full-stack de la aplicación web, desarrollando interfaces responsivas y servicios backend para procesamiento y análisis de datos ganaderos.}
\cvitem{Diseño e implementación de la arquitectura DevOps y containerización de la suite de servicios mediante Docker.}
\cvitem{Configuración y automatización de canalizaciones de integración y despliegue continuo (CI/CD) con GitHub Actions para despliegues automatizados en VPS.}

\jobsep

\cventry{Consultor DevOps y cloud}{Gipsycorp}{Venezuela · Remoto · Consultoría por proyecto}{Marzo 2026 – Julio 2026}
\cvitem{Migración integral de infraestructura y servicios corporativos desde Microsoft Azure hacia entorno VPS on-premise.}
\cvitem{Containerización y orquestación de la arquitectura de servicios mediante Docker.}
\cvitem{Automatización de canalizaciones de despliegue continuo (CD) mediante GitHub Actions hacia el VPS on-premise.}
\cvitem{Configuración y securización de la capa perimetral web con Azure Front Door para aceleración, protección WAF y balanceo de tráfico.}

\jobsep

\cventry{Auxiliar docente}{Universidad Central de Venezuela (UCV)}{Caracas, Venezuela · Presencial · Laboratorios docentes de computación}{Julio 2025 – Presente}
\cvitem{Dictado de preparadurías y orientación académica a estudiantes de la licenciatura en computación en la materia de Probabilidad y Estadística.}
\cvitem{Acompañamiento en la resolución de problemas analíticos, preparación de guías de ejercicios y apoyo a la cátedra docente en evaluaciones.}
\cvitem{Soporte técnico, configuración y mantenimiento de las computadoras de los laboratorios, asegurando su correcto funcionamiento para las clases y prácticas.}

\jobsep

\cventry{Operador telefónico JR de ATC}{Inversiones DRC S.A}{Caracas, Venezuela · Presencial}{Septiembre 2022 – Enero 2023}
\cvitem{Soporte y atención al cliente, gestión ágil de incidencias y cumplimiento de métricas de calidad de servicio.}

% ==============================================================================
% PÁGINA 2
% ==============================================================================
\newpage

% --- EDUCACIÓN ---
\cvsection{EDUCACIÓN}

\cventry{Licenciatura en computación}{Universidad Central de Venezuela (UCV)}{Caracas, Venezuela}{Enero 2023 – Presente}

% --- CURSOS ---
\cvsection{CURSOS}

\cvitem{\textbf{Machine learning y deep learning}: Análisis y preparación de datos, entrenamiento de modelos predictivos tradicionales y desarrollo de arquitecturas neuronales aplicadas (Facultad de Ciencias, UCV).}
\cvitem{\textbf{Bases de datos relacionales}: Diseño y modelado relacional, lenguaje SQL, normalización e integridad de datos (Facultad de Ciencias, UCV).}
\cvitem{\textbf{Programación con Python}: Creación de aplicaciones web interactivas con frameworks como Flask (abril 2022).}
\cvitem{\textbf{Git y GitHub}: Control de versiones, ramas, trabajo colaborativo y flujos de trabajo en equipo (julio 2021).}
\cvitem{\textbf{Desarrollo web}: Construcción de interfaces web dinámicas y adaptativas con HTML, CSS, JavaScript y Next.js (febrero 2021).}

% --- HABILIDADES TÉCNICAS ---
\cvsection{HABILIDADES TÉCNICAS}

\cvskill{Desarrollo web y móvil}{React.js, Next.js, React Native, Astro.js, Tailwind CSS, HTML5, CSS3.}
\cvskill{Backend y frameworks}{Express.js, NestJS, Django, Flask, Laravel, Node.js, APIs REST, GraphQL.}
\cvskill{DevOps e infraestructura}{Docker, Docker Compose, GitHub Actions (CI/CD), Microsoft Azure, Azure Front Door, VPS, Linux / Bash, Nginx.}
\cvskill{Datos e inteligencia artificial}{Machine learning, Deep learning, SQL, PostgreSQL, MySQL, SQLite.}
\cvskill{Lenguajes de programación}{JavaScript, TypeScript, Python, PHP, Java, C/C++, SQL.}
\cvskill{Herramientas y metodologías}{Git, GitHub, metodologías ágiles (Scrum).}

% --- IDIOMAS ---
\cvsection{IDIOMAS}

{\fontsize{9.6pt}{14.5pt}\selectfont Español (nativo)\cvdot Inglés (básico, lectura y comprensión técnica)}\par

\end{document}
